% Part: first-order-logic
% Chapter: sequent-calculus
% Section: proving-things-quant

\documentclass[../../../include/open-logic-section]{subfiles}

\begin{document}

\olfileid{fol}{seq}{prq}

\olsection{પરિમાણકો ધરાવતી \usetoken{P}{derivation}}

\begin{ex}
સિક્વન્ટ $\lexists[x][\lnot !A(x)] \Sequent
\lnot \lforall[x][!A(x)]$ની $\Log{LK}$-!!{derivation} આપો.

પરિમાણકો સાથે કામ કરતી વખતે આઇગનચલ-શરતનો ભંગ ન થાય તેની ખાતરી કરવી
પડે છે, અને ક્યારેક તે માટે અમુક અનુમાનો કયા ક્રમમાં કરવા તેની ગોઠવણ
બદલવી પડે છે. સામાન્ય રીતે આઇગનચલ-શરતને આધીન નિયમોને પહેલાં સંભાળવાનો
પ્રયત્ન કરવાથી મદદ મળે છે (પૂરી થયેલી સાબિતીમાં તેઓ વધુ નીચે હશે).
આરંભિક સિક્વન્ટ કેવું હશે તે આગળથી વિચારીને અનુમાન કરવાનો પ્રયત્ન કરવો
પણ સારો ઉપાય છે. આપણા કિસ્સામાં તે $!A(a) \Sequent !A(a)$ જેવું હોવું
પડશે. એટલે પરિમાણકના નિયમોને ``ઊલટા'' લાગુ કરતી વખતે $\lforall$ અને
$\lexists$ બંને નિયમ માટે એક જ પદ—જેને આપણે $a$ કહીશું—પસંદ કરવું પડશે.
દરેક નિયમ માટે જુદાં પદો પસંદ કરીએ તો અંતે $!A(a) \Sequent !A(b)$ જેવું
કંઈક મળે, જે અલબત્ત નિષ્પન્ન થઈ શકતું નથી.

હંમેશની જેમ આ લખીને શરૂ કરીએ:
\begin{prooftree}
\AxiomC{}
\UnaryInf$\lexists[x][\lnot !A(x)] \fCenter \lnot \lforall[x][!A(x)]$
\end{prooftree}
અહીં \LeftR{\exists} અથવા \RightR{\lnot}, બંનેમાંથી ગમે તે નિયમ લગાડી
શકીએ. \LeftR{\exists} નિયમ આઇગનચલ-શરતને આધીન છે, તેથી તેને મોડું કરવાને
બદલે વહેલું સંભાળવું સારો ઉપાય છે; એટલે પહેલાં એ જ નિયમ લગાડીએ.
\begin{prooftree}
\AxiomC{}
\UnaryInf$ \lnot !A(a) \fCenter \lnot \lforall[x][!A(x)]$
\RightLabel{\LeftR{\lexists}}
\UnaryInf$ \lexists[x][\lnot !A(x)] \fCenter \lnot \lforall[x][!A(x)]$
\end{prooftree}
\LeftR{\lnot} અને \RightR{\lnot} નિયમોને ઊલટી દિશામાં લગાડતાં મળે છે:
\begin{prooftree}
\AxiomC{}
\UnaryInf$\lforall[x][!A(x)] \fCenter !A(a)$
\RightLabel{\LeftR{\lnot}}
\UnaryInf$\lnot !A(a), \lforall[x][!A(x)] \fCenter $
\RightLabel{\LeftR{\Exchange}}
\UnaryInf$\lforall[x][!A(x)], \lnot !A(a) \fCenter $
\RightLabel{\RightR{\lnot}}
\UnaryInf$ \lnot !A(a) \fCenter \lnot \lforall[x] !A(x)$
\RightLabel{\LeftR{\lexists}}
\UnaryInf$ \lexists[x] \lnot !A(x) \fCenter \lnot \lforall[x] !A(x)$
\end{prooftree}
હવે આપણી એકમાત્ર શક્યતા \LeftR{\forall} નિયમ લગાડવાની છે. આ નિયમ
આઇગનચલ-પ્રતિબંધને આધીન નથી, તેથી અહીં કોઈ અડચણ નથી. યાદ રાખો કે આપણે
$!A(a) \Sequent !A(a)$ સ્વરૂપનું આરંભિક સિક્વન્ટ મેળવવા માગીએ છીએ;
તેથી નિયમ લગાડીએ ત્યારે $!A$ના પદ તરીકે $a$ પસંદ કરવું જોઈએ.
\begin{prooftree}
\Axiom$!A(a) \fCenter !A(a)$
\RightLabel{\LeftR{\lforall}}
\UnaryInf$\lforall[x][!A(x)] \fCenter !A(a)$
\RightLabel{\LeftR{\lnot}}
\UnaryInf$\lnot !A(a), \lforall[x][!A(x)] \fCenter $
\RightLabel{\LeftR{\Exchange}}
\UnaryInf$\lforall[x][!A(x)], \lnot !A(a) \fCenter $
\RightLabel{\RightR{\lnot}}
\UnaryInf$ \lnot !A(a) \fCenter \lnot \lforall[x][!A(x)]$
\RightLabel{\LeftR{\lexists}}
\UnaryInf$ \lexists[x][ \lnot !A(x)] \fCenter \lnot \lforall[x][!A(x)]$
\end{prooftree}
ખાસ કરીને પરિમાણકો સાથે કામ કરતી વખતે, હવે આઇગનચલ-શરતનો ભંગ થયો નથી
તે ફરી તપાસવું જરૂરી છે. આપણે લગાડેલો આઇગનચલ-શરતને આધીન એકમાત્ર નિયમ
\LeftR{\exists} હતો, અને આઇગનચલ~$a$ તેના નીચલા સિક્વન્ટમાં
(અંતિમ સિક્વન્ટમાં) આવતું નથી. તેથી આ યોગ્ય !!{derivation} છે.
\end{ex}

\begin{prob}
નીચેના સિક્વન્ટોની !!{derivation}s આપો:
\begin{enumerate}
\item $\Sequent (\lforall[x][!A(x)] \land \lforall[y][!B(y)]) \lif
\lforall[z][(!A(z) \land !B(z))]$.
\item $\Sequent (\lexists[x][!A(x)] \lor \lexists[y][!B(y)]) \lif
\lexists[z][(!A(z) \lor !B(z))]$.
\item $\lforall[x][(!A(x) \lif !B)] \Sequent \lexists[y][!A(y)] \lif !B$.
\item $\lforall[x][\lnot !A(x)] \Sequent \lnot\lexists[x][!A(x)]$.
\item $\Sequent \lnot\lexists[x][!A(x)] \lif \lforall[x][\lnot !A(x)]$.
\item $\Sequent \lnot\lexists[x][\lforall[y][((!A(x,y) \lif \lnot
!A(y,y)) \land (\lnot !A(y,y) \lif !A(x,y)))]]$.
\end{enumerate}
\end{prob}

\begin{prob}
નીચેના સિક્વન્ટોની !!{derivation}s આપો:
\begin{enumerate}
\item $\Sequent \lnot\lforall[x][!A(x)] \lif \lexists[x][\lnot!A(x)]$.
\item $(\lforall[x][!A(x)] \lif !B) \Sequent \lexists[y][(!A(y) \lif !B)]$.
\item $\Sequent \lexists[x][(!A(x) \lif \lforall[y][!A(y)])]$.
\end{enumerate}
(આ બધામાં $\RightR{\Contraction}$ નિયમ જરૂરી છે.)
\end{prob}

\end{document}
